\section{Conclusion}
\label{sec:conclusion}

Training with a learned physical reward improves measured utility per RTL
draw on the studied streaming-accelerator tasks under the sampled oracle
contract. On the sealed split, the RF package increases failure-penalized
post-route timing-derived frequency from
\studyclaim{StudyTwoOverallSftPenalized} to
\studyclaim{StudyTwoOverallRfPenalized}~MHz, with positive effects for both
policy-training seeds. The gain comes from faster oracle-passing
implementations, despite a slightly lower pass rate. The original-MLP and
single-seed correctness-only controls help explain this result without
isolating the individual components of the RF repair. A separate setup-closure
extension repeats the positive RF and negative correctness-only effects in
two additional seeds per arm, all sharing the same SFT initialization. The
Qwen-Coder study provides further evidence that the overall training recipe
can benefit another backbone.

The benefits vary across families and come with implementation costs,
including higher LUT and register use, slightly higher vectorless power, and
a small regression in the matched VerilogEval experiment. The setup-closure
extension and descriptive board measurements provide additional physical
evidence while retaining their separate contracts and limits; neither replaces
the sealed endpoint or establishes full timing sign-off. Within these bounds,
the results support using a learned physical reward to improve the sampling
distribution, with executable testing, complete failure accounting, and
independent physical measurements determining the reported outcome.
